\documentclass[11pt]{article}
\usepackage[margin=1in]{geometry}
\usepackage{amsmath,amssymb,amsthm}
\usepackage{booktabs}
\usepackage[round,authoryear]{natbib}
\usepackage{xcolor}
\usepackage[hidelinks]{hyperref}

\newtheorem{proposition}{Proposition}
\newtheorem{corollary}{Corollary}
\theoremstyle{remark}
\newtheorem{remark}{Remark}

% Open items for the authors and the desks; every \todo is listed in F23_HANDOFF.md.
\newcommand{\todo}[1]{\textcolor{red}{[\textbf{TODO:} #1]}}
\newcommand{\notewidth}{0.95\textwidth} % width of table notes

% output/macros.tex defines four names containing the digit 0 (\PwHomoH0Falla, \PwHetH0Falla, \PwHomoPowH0, \PwHetPowH0),
% which LaTeX cannot parse as control words. Until code/analysis/02_paper_macros.py renames them (see F23_HANDOFF.md), the
% file is read with 0 as a letter and those four macros are aliased. No value is changed.
\catcode`\0=11
% Generado por code/analysis/02_paper_macros.py. No editar a mano. Fuentes y nombres: ver ese script.
\newcommand{\bLinEst}{\ensuremath{-1.655}}
\newcommand{\bLinSe}{1.786}
\newcommand{\bLinDof}{114}
\newcommand{\bLinP}{0.178}
\newcommand{\bLinCiLo}{\ensuremath{-5.193}}
\newcommand{\bLinCiHi}{1.883}
\newcommand{\bCLiqEst}{\ensuremath{-2.874}}
\newcommand{\bCLiqSe}{1.653}
\newcommand{\bCLiqDof}{112}
\newcommand{\bCLiqP}{0.042}
\newcommand{\bCLiqCiLo}{\ensuremath{-6.149}}
\newcommand{\bCLiqCiHi}{0.400}
\newcommand{\bULiqEst}{0.841}
\newcommand{\bULiqSe}{4.341}
\newcommand{\bULiqDof}{112}
\newcommand{\bULiqP}{0.577}
\newcommand{\bULiqCiLo}{\ensuremath{-7.760}}
\newcommand{\bULiqCiHi}{9.441}
\newcommand{\bCLiqNovEst}{\ensuremath{-4.981}}
\newcommand{\bCLiqNovSe}{2.814}
\newcommand{\bCLiqNovDof}{99}
\newcommand{\bCLiqNovP}{0.040}
\newcommand{\bCLiqNovCiLo}{\ensuremath{-10.565}}
\newcommand{\bCLiqNovCiHi}{0.603}
\newcommand{\bULiqNovEst}{1.024}
\newcommand{\bULiqNovSe}{4.687}
\newcommand{\bULiqNovDof}{99}
\newcommand{\bULiqNovP}{0.586}
\newcommand{\bULiqNovCiLo}{\ensuremath{-8.276}}
\newcommand{\bULiqNovCiHi}{10.324}
\newcommand{\bCRobFromEst}{\ensuremath{-2.035}}
\newcommand{\bCRobFromSe}{1.974}
\newcommand{\bCRobFromDof}{102}
\newcommand{\bCRobFromP}{0.152}
\newcommand{\bCRobFromCiLo}{\ensuremath{-5.950}}
\newcommand{\bCRobFromCiHi}{1.879}
\newcommand{\bURobFromEst}{0.059}
\newcommand{\bURobFromSe}{4.097}
\newcommand{\bURobFromDof}{102}
\newcommand{\bURobFromP}{0.506}
\newcommand{\bURobFromCiLo}{\ensuremath{-8.067}}
\newcommand{\bURobFromCiHi}{8.184}
\newcommand{\bCRobQuartEst}{\ensuremath{-6.085}}
\newcommand{\bCRobQuartSe}{3.045}
\newcommand{\bCRobQuartDof}{112}
\newcommand{\bCRobQuartP}{0.024}
\newcommand{\bCRobQuartCiLo}{\ensuremath{-12.118}}
\newcommand{\bCRobQuartCiHi}{\ensuremath{-0.052}}
\newcommand{\bURobQuartEst}{\ensuremath{-0.840}}
\newcommand{\bURobQuartSe}{2.488}
\newcommand{\bURobQuartDof}{112}
\newcommand{\bURobQuartP}{0.368}
\newcommand{\bURobQuartCiLo}{\ensuremath{-5.769}}
\newcommand{\bURobQuartCiHi}{4.089}
\newcommand{\bCRobLogitEst}{\ensuremath{-4.576}}
\newcommand{\bCRobLogitSe}{2.196}
\newcommand{\bCRobLogitDof}{112}
\newcommand{\bCRobLogitP}{0.020}
\newcommand{\bCRobLogitCiLo}{\ensuremath{-8.928}}
\newcommand{\bCRobLogitCiHi}{\ensuremath{-0.225}}
\newcommand{\bURobLogitEst}{1.890}
\newcommand{\bURobLogitSe}{5.289}
\newcommand{\bURobLogitDof}{112}
\newcommand{\bURobLogitP}{0.639}
\newcommand{\bURobLogitCiLo}{\ensuremath{-8.589}}
\newcommand{\bURobLogitCiHi}{12.369}
\newcommand{\bCRobLagsEst}{\ensuremath{-2.219}}
\newcommand{\bCRobLagsSe}{1.582}
\newcommand{\bCRobLagsDof}{104}
\newcommand{\bCRobLagsP}{0.082}
\newcommand{\bCRobLagsCiLo}{\ensuremath{-5.356}}
\newcommand{\bCRobLagsCiHi}{0.917}
\newcommand{\bURobLagsEst}{3.034}
\newcommand{\bURobLagsSe}{4.028}
\newcommand{\bURobLagsDof}{104}
\newcommand{\bURobLagsP}{0.773}
\newcommand{\bURobLagsCiLo}{\ensuremath{-4.953}}
\newcommand{\bURobLagsCiHi}{11.021}
\newcommand{\verdictoVs}{h0\_falla}
\newcommand{\DiagFullN}{146}
\newcommand{\DiagFullDoneRestr}{73}
\newcommand{\DiagFullDoneHolg}{73}
\newcommand{\DiagFullDtwoRestr}{31.036}
\newcommand{\DiagFullDtwoHolg}{21.754}
\newcommand{\DiagFullDthreeF}{0.50}
\newcommand{\DiagFullDthreeP}{0.606}
\newcommand{\DiagFullLB}{16.78}
\newcommand{\DiagFullLBP}{0.158}
\newcommand{\DiagFullPredF}{1.07}
\newcommand{\DiagFullPredDfA}{12}
\newcommand{\DiagFullPredDfB}{133}
\newcommand{\DiagFullPredP}{0.391}
\newcommand{\NShockFull}{146}
\newcommand{\DiagNoVN}{133}
\newcommand{\DiagNoVDoneRestr}{62}
\newcommand{\DiagNoVDoneHolg}{71}
\newcommand{\DiagNoVDtwoRestr}{18.195}
\newcommand{\DiagNoVDtwoHolg}{19.810}
\newcommand{\DiagNoVDthreeF}{0.49}
\newcommand{\DiagNoVDthreeP}{0.613}
\newcommand{\DiagNoVLB}{22.86}
\newcommand{\DiagNoVLBP}{0.029}
\newcommand{\DiagNoVPredF}{1.66}
\newcommand{\DiagNoVPredDfA}{12}
\newcommand{\DiagNoVPredDfB}{120}
\newcommand{\DiagNoVPredP}{0.085}
\newcommand{\NShockNoV}{133}
\newcommand{\sdEps}{0.6001}
\newcommand{\betaRef}{\ensuremath{-1.820}}
\newcommand{\betaRefSd}{\ensuremath{-1.092}}
\newcommand{\piZero}{0.949}
\newcommand{\piZeroSe}{0.248}
\newcommand{\PwHomoHzeroFalla}{0.824}
\newcommand{\PwHomoConfirma}{0.037}
\newcommand{\PwHomoNoConc}{0.139}
\newcommand{\PwHomoPowHzero}{0.176}
\newcommand{\PwHomoPowHone}{0.118}
\newcommand{\mdeEpsUnit}{1.82}
\newcommand{\mdeFfrUnit}{1.92}
\newcommand{\PwHetHzeroFalla}{0.854}
\newcommand{\PwHetConfirma}{0.059}
\newcommand{\PwHetNoConc}{0.087}
\newcommand{\PwHetPowHzero}{0.146}
\newcommand{\PwHetPowHone}{0.121}
\newcommand{\MdeHomoC}{5.5}
\newcommand{\MdeHomoEps}{10.01}
\newcommand{\MdeHomoFfr}{10.54}
\newcommand{\MdeHomoRivalMax}{0.035}
\newcommand{\MdeHetC}{5.2}
\newcommand{\MdeHetEps}{9.46}
\newcommand{\MdeHetFfr}{9.97}
\newcommand{\MdeHetRivalMax}{0.033}
\newcommand{\SizeHomo}{0.004}
\newcommand{\SizeHet}{0.011}
\newcommand{\CovAumNwHcLo}{0.948}
\newcommand{\CovAumNwHcHi}{0.990}

\let\PwHomoHzeroFalla\PwHomoH0Falla
\let\PwHetHzeroFalla\PwHetH0Falla
\let\PwHomoPowHzero\PwHomoPowH0
\let\PwHetPowHzero\PwHetPowH0
\catcode`\0=12

\title{Precautionary Reserves and Monetary Transmission to Financially Constrained Firms:\\
Theory and a Preregistered, Bounded Measurement}
\author{\todo{author list and affiliations}}
\date{Working paper draft}

\begin{document}
\maketitle

\begin{abstract}
\noindent
In models where a binding financing constraint fixes investment, monetary policy reaches constrained firms only through
the repricing of their existing debt. We show that this prediction is a corner. When firms hold precautionary reserves that
they can adjust, the investment response of a constrained firm is a weighted average of the unconstrained response and the
repricing channel, and the weight on the unconstrained response is bounded below by a strictly positive constant that
depends only on primitives. Aggregating firms that borrow from a common bank adds a crowding-out term that dampens the aggregate
response but does not reverse it when repricing does not hurt constrained firms. All results are verified symbolically. We then report a preregistered measurement for
the United States, 1969--2005: the average response of non-residential investment, four to twelve quarters ahead, to
narrative monetary policy shocks, overall and by the state of corporate liquidity. The preregistered gate fails: the
linear response is negative but its confidence interval contains both zero and the reference effect. Under the
preregistered rule the state-dependent hypothesis is therefore not evaluated, and we report its estimates as measurement
without a verdict. With the power computed before estimation, this was the most likely outcome even if the reference effect
were true. The design had 80 percent power only for responses of roughly 10 percentage points per percentage point of
shock, about five times the reference effect.
\end{abstract}

\section{Introduction}\label{sec:intro}

How much of a monetary policy impulse reaches firms whose investment is limited by a financing constraint? Two polar answers
are common. In models where the constraint pins investment down, an interest rate change moves the investment of
constrained firms only through the cash flow freed or absorbed by the repricing of their outstanding debt; in the floating-rate
channel of \citet[WP pp.~10--11]{IppolitoOzdagliPerezOrive2018}, the investment of a constrained firm moves one for one
with the change in its interest bill. Without a constraint, investment responds through the usual cost-of-capital margin.
Which of the two describes the economy is then a question about how many firms are constrained.

This paper makes two contributions. The first is theoretical. We study a firm that faces a binding financing capacity and
holds a precautionary liquidity reserve that it can adjust. Its investment response to a monetary easing is a weighted average
of the unconstrained response and the repricing channel
(Proposition~\ref{prop:transmission}). The weight on the unconstrained response falls as reserves become costlier to run down,
but it is bounded below by a strictly positive constant that depends only on primitives (Proposition~\ref{prop:limits}). The
rigid corner, in which a constrained firm without floating-rate debt does not respond at all, is reached only in the limit or at
a corner of the reserve choice, never at an interior optimum. Aggregating heterogeneous firms that borrow from a common bank
adds a crowding-out term that dampens the aggregate response but keeps it positive whenever repricing does not hurt
constrained firms (Proposition~\ref{prop:aggregation}). Every derivation is verified symbolically, and the aggregation formula
is also checked against a numerically solved nonlinear equilibrium (Appendix~\ref{app:verification}).

The second contribution is a measurement. The theory implies that a monetary tightening lowers aggregate investment and that it
lowers investment also when many firms are constrained. We take both implications to U.S. data with narrative monetary policy
shocks \citep{RomerRomer2004}, using a design that was preregistered, with its decision rule and its power, before the
investment response was estimated. The statistic is the average response of non-residential investment over horizons of four
to twelve quarters. The state is whether the liquid assets of non-financial corporations relative to their debt are below
their recent trend. The decision rule first requires the linear response to be negative (a gate) and only then evaluates the
response in the low-liquidity state, which must hold both in the full sample and without the Volcker disinflation.

The gate fails. The linear estimate is \bLinEst{} percentage points (pp) of average investment per pp of shock, with a
one-sided $p$-value of \bLinP{} and a 95 percent confidence interval of $[\bLinCiLo, \bLinCiHi]$, which contains both zero and the
reference effect of \betaRef{} that we fixed ex ante from \citet{TenreyroThwaites2016}. The hypothesis about the constrained state is
therefore not interpreted. We report all six preregistered estimates of the response in the low-liquidity state, with
their $p$-values and intervals, as measurement without a verdict. This outcome is not evidence that monetary policy does not
affect investment. Before estimation, our Monte Carlo put the probability of exactly this outcome at \PwHomoHzeroFalla{} if the
true response were the reference effect (\PwHetHzeroFalla{} with heteroskedastic noise). The design confirms H1 with probability
0.8 only if the responses are of roughly 10 pp per pp of shock, about five times the reference effect. % num-ok: 0.8 es la potencia objetivo que define el MDE en el prerregistro, no un resultado

We see the measurement as useful for two reasons. It bounds what aggregate U.S. data on this question can say, and it shows,
with a design fixed in advance, why a sample of this length cannot distinguish the polar answers above. Two extensions that could
add power, a cross-section of industries by firm size and a panel of euro-area countries with contractual measures of
repricing exposure, did not pass the data requirements that we set for them before looking at any outcome
(Section~\ref{sec:conclusion}).

\paragraph{Related literature.}
A large literature studies how firm balance sheets and debt structure shape the investment response to monetary policy
\citep{OttonelloWinberry2020,CloyneFerreiraFroemelSurico2023,GrobZullig2024,JungherrMeierReineltSchott2024,GreenwaldKrainerPaul2025,BotteroEtAl2026}.
\citet{OttonelloWinberry2020} find that firms with low default risk are the most responsive. \todo{Jeenas (BSE WP 1409) on
firm liquidity and the investment response is the closest empirical paper; the bib entry is not VERIFIED (publication status
in dispute) and cannot be cited yet.} Our model nests the floating-rate channel of
\citet{IppolitoOzdagliPerezOrive2018}: their constrained firm, whose capital equals net worth plus a fixed borrowing limit
(WP eqs.~2 and 7, pp.~10--11), is the rigid limit of ours with fully floating debt. \citet{PerezOriveTimmerVanDerGhote2026}
show that when several financing constraints bind, firms respond more to tightenings than to easings, and that with a single
binding constraint the response is symmetric (Propositions 1 and 2, WP pp.~12--13). Our model has a single capacity
constraint and an interior reserve margin; we pool tightenings and easings and leave asymmetry to exploratory work. On the
empirical side, we build on state-dependent local projections with narrative shocks \citep{TenreyroThwaites2016,Ramey2016}
and on the use of future shocks as controls proposed by \citet{TeulingsZubanov2014}. \todo{cite the authors' earlier work (TFG1,
TFG2) on which the model builds, once there are verified bib entries.}

\section{Model}\label{sec:model}

\subsection{A constrained firm with a reserve margin}

A firm chooses investment $I$ and a precautionary liquidity reserve $m$. Each unit of reserves uses $\gamma>0$ units of
financing capacity $H$, so that $I+\gamma m\le H$. The marginal return on investment is $q$ and its marginal cost rises with
slope $a=b+\eta v+\psi/N$, where $b$ is the curvature of the project cost, $\eta v$ an uncertainty term, and $\psi/N$ the slope
of the lending rate charged by a common bank (Section~\ref{sec:aggregation}). Reserves insure against a liquidity need
$\xi\sim N(\mu,\sigma^2)$. Let $z=(m-\mu)/\sigma$, let $Q(z)$ be the upper tail of the standard normal and let
$G(m)=\mathbb{E}[(\xi-m)^+]$ be the expected shortfall, so that $G'(m)=-Q(z)$. With the capacity constraint binding
($\lambda>0$) and an interior reserve, the optimum satisfies
\begin{equation}\label{eq:kkt}
\lambda = q-aI,\qquad \kappa_F\,G(m)=\delta+\gamma\lambda,\qquad I+\gamma m=H,
\end{equation}
where $\delta$ is the carrying cost of reserves and $\kappa_F>0$ scales the value of insurance. Define
$B\equiv-\kappa_F G'(m)=\kappa_F Q(z)$, the curvature of the reserve value at the optimum.

A monetary easing is a small change $u$ in the policy rate (a cut). It raises the net return on investment, $q_u=p_f>0$, and it
relaxes financing capacity through the repricing of debt, $H_u=g$, where $g=\gamma\mathcal{E}_D$ and $\mathcal{E}_D$ is the
exposure of the firm's interest bill to the policy rate. We hold $a$, $\sigma$ and $\delta$ fixed.

\begin{proposition}[Transmission with a reserve margin]\label{prop:transmission}
Differentiating \eqref{eq:kkt} with respect to $u$,
\begin{equation}\label{eq:Iu}
I_u=\frac{\gamma^2 p_f+B g}{B+\gamma^2 a}=\omega\,\frac{p_f}{a}+(1-\omega)\,g,\qquad
\omega=\frac{\gamma^2 a}{B+\gamma^2 a},
\end{equation}
and $m_u=\gamma(a g-p_f)/(B+\gamma^2 a)$, $\lambda_u=B(p_f-a g)/(B+\gamma^2 a)$.
\end{proposition}

The response is a weighted average of two margins: the price margin $p_f/a$, which is the response of an unconstrained firm, and
the cash-flow margin $g$. The weight $\omega$ is high when reserves are cheap to adjust ($B$ small): the firm then funds
investment by running down reserves and behaves as if unconstrained.

\begin{proposition}[Limits and a lower bound on the weight]\label{prop:limits}
\begin{enumerate}
\item[(i)] As $B\to0$, $I_u\to p_f/a$: the unconstrained response.
\item[(ii)] As $B\to\infty$, $I_u\to g$. With $g=0$ this is the zero response of models in which the constraint fixes
investment. With $\gamma=1$ and $\mathcal{E}_D=L_0$ (all debt $L_0$ floating), it is the response $\partial K_1/\partial r_1=-L_0$
implied by equations 2 and 7 of \citet[WP pp.~10--11]{IppolitoOzdagliPerezOrive2018}. With an adjustable reserve and finite
$B$, their response is not exact.
\item[(iii)] Because $Q$ decreases strictly from 1 to 0, $B=\kappa_F Q(z)\in(0,\kappa_F)$ and therefore
\begin{equation}\label{eq:wmin}
\omega\in(\omega_{\min},1),\qquad \omega_{\min}=\frac{\gamma^2 a}{\kappa_F+\gamma^2 a}>0 .
\end{equation}
The rigid limit $\omega=0$ is not reached at an interior optimum with finite $\kappa_F$. It obtains exactly at a corner where
the reserve is fixed, in which case $I_u=g$.
\end{enumerate}
\end{proposition}

\begin{corollary}[Exposure gradient]\label{cor:gradient}
$\partial I_u/\partial g=1-\omega$, which is strictly positive when $B>0$ and zero in the limit $B\to0$.
\end{corollary}

\begin{remark}
With $g\ge0$, \eqref{eq:Iu} and \eqref{eq:wmin} give $I_u\ge\omega\,p_f/a>\omega_{\min}\,p_f/a>0$: at an interior optimum a
constrained firm always transmits a positive fraction of the unconstrained response.
\end{remark}

\subsection{Aggregation with a common bank}\label{sec:aggregation}

Firms borrow from a common bank at $r_L=r_f+\phi L$, where $L$ is aggregate lending and $\phi=\psi/N$. Each firm takes $r_L$ as
given, so its own cost curvature is $a_j$ (without $\phi$), and $\phi$ enters only through the price. There are $K$ types with
mass $s_j$, curvature $a_j$, exposure $g_j$ and weight $\omega_j=\gamma_j^2 a_j/(B_j+\gamma_j^2 a_j)$, with $\omega_j=1$ for a
type whose constraint is slack. Define $\Theta=\sum_j s_j\,\omega_j/a_j$ and $\bar G=\sum_j s_j(1-\omega_j)g_j$.

\begin{proposition}[Aggregate response]\label{prop:aggregation}
\begin{equation}\label{eq:Lu}
L_u=\frac{\bar G+\Theta\,p_f}{1+\phi\,\Theta},
\end{equation}
and each type responds as in \eqref{eq:Iu} with $p_f$ replaced by the effective price shock $x=p_f-\phi L_u$. Moreover:
\begin{enumerate}
\item[(i)] With $g_j=0$ for all $j$, $L_u<\Theta\,p_f$: the common bank dampens the aggregate response by the factor
$1/(1+\phi\Theta)$.
\item[(ii)] With a single type, \eqref{eq:Lu} coincides with \eqref{eq:Iu} for a representative firm with curvature
$a_j+\phi$.
\item[(iii)] If $g_j\ge0$ and $\omega_j\ge\omega_{\min}>0$ for all $j$, then $L_u>0$.
\end{enumerate}
\end{proposition}

The crowding-out term matters for the measurement. When constrained firms receive cash-flow relief, their extra borrowing raises
the lending rate and lowers the investment of slack firms; the aggregate response is smaller than the sum of individual
responses at a fixed lending rate. Proposition~\ref{prop:aggregation} also corrects a shortcut: the aggregate formula that treats
constrained firms as rigid ($\omega_j=0$) does not hold when reserves are adjustable. In a nonlinear equilibrium solved
numerically, \eqref{eq:Lu} matches the finite-difference derivative and the rigid formula does not
(Appendix~\ref{app:verification}).

\paragraph{Scope.} The results are local: the mass of constrained firms is held fixed, although in the data it moves with the
shock. They are partial equilibrium, with the price of investment fixed, and $B_j$ is evaluated at the initial equilibrium.

\subsection{From the model to the hypotheses}\label{sec:hypotheses}

A monetary tightening is $-u$. Our sign convention for the estimated responses is that $\beta<0$ means investment falls after a
tightening and rises after an easing. Proposition~\ref{prop:aggregation}(iii) gives $L_u>0$ whenever repricing does not hurt
constrained firms ($g_j\ge0$) and reserves are interior, which yields
\begin{itemize}
\item[\textbf{H0}] the average response of aggregate investment to a monetary shock is negative, $\beta_{\text{lin}}<0$.
\end{itemize}
The remark after Corollary~\ref{cor:gradient} gives $I_u>0$ for a constrained firm with $g\ge0$, whereas the rigid corner predicts
$I_u=g$, which is zero without floating-rate debt. This yields
\begin{itemize}
\item[\textbf{H1}] in the state where many firms are constrained, the response is also negative, $\beta_C<0$.
\end{itemize}
H1 rests on the maintained assumption that net exposure to repricing is non-negative, $g\ge0$; a rejection would reject the model
with $g\ge0$, not the reserve margin as such. The theory does not sign the difference between the constrained and the unconstrained
state, which depends on whether $g$ is larger or smaller than $p_f/a$; that difference is therefore not a preregistered hypothesis.
The empirical state below measures aggregate corporate liquidity, which we use as a proxy for a large mass of firms on the
constrained margin. That mapping is an assumption of the measurement, not a result of the model.

\section{Preregistered design}\label{sec:design}

The hypotheses, the specification, the decision rule, the power and the minimum detectable effect were written and frozen
before the investment response was estimated (preregistration frozen at commit \texttt{3d6b508}, approved by the authors on
2026-11-08). The design replaces an earlier preregistration, based on high-frequency shocks \citep{BauerSwanson2023} and
short horizons, which was abandoned before any estimation because its Monte Carlo power against the reference effect was close
to its size. All deviations are dated and documented in the project specification.

\subsection{Data and sample}

The shock $\varepsilon_t$ is the \citet{RomerRomer2004} narrative series, defined as the residual from regressing the intended
change in the federal funds rate at each FOMC meeting on Greenbook forecasts, in the version extended through 2007 that is
distributed with the replication files of \citet{Ramey2016}. \todo{verified citation for the provenance of the extended series
(Wieland's update, openICPSR 135741), which the Literature desk could not open.} The quarterly shock is the sum of the
monthly shocks in the quarter, so its unit is the percentage point of federal funds summed over the quarter. Investment is the
BEA quantity index of real private non-residential fixed investment; controls are real GDP, the GDP deflator and the quarterly
average federal funds rate (FRED). Corporate liquidity is the ratio of broad liquid assets to credit-market debt of non-financial
corporations in the Financial Accounts of the United States (Z.1). \todo{formal citations for FRED, BEA and Z.1; no verified bib
entries exist yet.} All series are in their latest vintage, not real time.

The estimation sample is $t\in$ 1969Q1--2005Q2, the last quarter for which $t+12$ ends before 2008Q3; it has \NShockFull{}
quarters, \DiagFullDoneRestr{} in each state. Because most Volcker-era quarters fall in the low-liquidity state and carry a
large share of the shock variation (Table~\ref{tab:diagnostics}), the design also uses a sample without 1979Q4--1982Q4, with
\DiagNoVN{} quarters. The sensitivity of Romer--Romer-based estimates to the 1979--1982 period is discussed by
\citet{Coibion2012}. Dropping the Volcker quarters only removes rows; the state, the shock leads and lags and the dependent
variable of the remaining rows are computed from the full series. Missing shocks (before 1969 in the lags of the first rows,
2008 in the leads of the last two) are set to zero.

\subsection{State}

Let $s_t$ be the log of the liquidity ratio minus its average over the previous twelve quarters. The low-liquidity
(``constrained'') state is $D_{t-1}=1$ if $s_{t-1}$ is below its median in the full shock sample; the median is not recomputed
without Volcker. The state is predetermined with respect to the shock, but it is not randomly assigned: a difference between
states describes how the response to a shock varies with liquidity, not the effect of liquidity.

\subsection{Specification}

The dependent variable is the average response over horizons four to twelve,
\[
\bar y_t=\frac{1}{9}\sum_{h=4}^{12}100\left[\log I_{t+h}-\log I_{t-1}\right].
\]
The linear specification (\texttt{s1r\_lin}) is
\[
\bar y_t=\alpha+\beta_{\text{lin}}\varepsilon_t+\sum_{j=1}^{10}\lambda_j\varepsilon_{t-j}+\sum_{j=1}^{12}\varphi_j\varepsilon_{t+j}
+\sum_{\ell=1}^{2}\Gamma_\ell' x_{t-\ell}+u_t,
\]
and the state-dependent specification (\texttt{s1r\_liq}) replaces $\beta_{\text{lin}}\varepsilon_t$ with
$\delta D_{t-1}+\beta_C D_{t-1}\varepsilon_t+\beta_U(1-D_{t-1})\varepsilon_t$; \texttt{s1r\_liq\_nov} is the same specification
without the Volcker quarters. The controls $x_t$ are the log changes ($\times100$) of investment, real GDP and the deflator and the
change in the federal funds rate. The shock lags and leads are not interacted with the state. Estimation is by OLS.

\paragraph{Estimand and identifying assumption.}
Because the regression controls for the shocks of $t+1,\dots,t+12$, the coefficients on $\varepsilon_t$ measure the average
response over horizons four to twelve to an \emph{isolated} shock, holding future shocks fixed. This equals the total response
(that of a local projection without leads) only if the shocks are serially uncorrelated in the population, which is our declared
identifying assumption. Without the leads, the window projection is biased by the sample autocorrelation of the shocks. The idea
of controlling for intervening shocks is that of \citet{TeulingsZubanov2014}, who prove it for independent shocks in a panel with
fixed effects; we use it as motivation, not as a proof that the estimator is unbiased when shocks are autocorrelated.

\paragraph{Inference.}
Standard errors are Newey--West with a Bartlett kernel and five lags in calendar quarters, applied to HC3-rescaled residuals,
with Student-$t$ critical values on $N-k$ degrees of freedom. \todo{Newey and West (1987): bib entry not VERIFIED.} In our Monte Carlo
these intervals cover close to or above the nominal level for the coefficients that decide, without material undercoverage. We
accept the overcoverage because it costs power, not size. \todo{macro for the simulated coverage range of NW5+HC3 (F9), if the text should
give it.}

\paragraph{Shock diagnostics.}
Without using investment, we report a Ljung--Box test with twelve lags on $\varepsilon_t$; a predictability test, the Wald $F$ of
regressing $\varepsilon_t$ on four lags of real GDP growth, inflation and the change in the federal funds rate; the number of
quarters and the sum of $|\varepsilon_t|$ by state; and a balance test of $\varepsilon_t$ on $D_{t-1}$ and $s_{t-1}$. It was
fixed in advance that a rejection would change neither the specification nor the decision rule. It only qualifies the
comparison of an isolated-shock estimate with the reference effect, which is a total response.

\subsection{Reference effect}

The reference effect translates a published estimate into our units. \citet[WP pp.~11--12]{TenreyroThwaites2016} report, in their linear model, a fall of
2.5 percent in household durable spending and business investment after a 1 pp rise in the federal funds rate, % num-ok: cifra citada de Tenreyro-Thwaites (2016), no es un resultado de este paper
with the shock scaled to raise the funds rate by 1 pp on impact; their text does not state the horizon at which that figure is
read. We multiply it by the average over $h=4,\dots,12$ of a linear ramp that peaks at $h=10$ (an assumption of ours) and by the
impact response of the funds rate to our shock, $\pi_0=\piZero$ (s.e.\ \piZeroSe). The result is
$\beta_{\text{ref}}=\betaRef$ pp of $\bar y$ per pp of shock; a one-standard-deviation shock ($\sigma_\varepsilon=\sdEps$) then
moves $\bar y$ by \betaRefSd{} pp. A multiple $c$ of the reference effect is written $c\,\beta_{\text{ref}}$.

\subsection{Decision rule}

The rule fixes a one-sided level of 0.05 with $t$ critical values and a single H1, without multiplicity correction: % num-ok: nivel prerregistrado del contraste
\begin{enumerate}
\item \textbf{H0 gate.} If $\hat\beta_{\text{lin}}$ in \texttt{s1r\_lin} (full sample) is not negative with $p<0.05$, the outcome % num-ok: nivel prerregistrado del contraste
is ``H0 is not confirmed'': H1 is not interpreted, and $\hat\beta_C$, $\hat\beta_U$ and their intervals are published as
measurement without a verdict.
\item Otherwise, $\hat\beta_C$ in \texttt{s1r\_liq} is evaluated: if it is positive with $p<0.05$, H1 is refuted; if it is % num-ok: nivel prerregistrado del contraste
negative with $p<0.05$, the rule moves to step 3; otherwise the outcome is inconclusive. % num-ok: nivel prerregistrado del contraste
\item H1 is confirmed only if $\hat\beta_C$ is also negative with $p<0.05$ in \texttt{s1r\_liq\_nov}; otherwise the outcome is % num-ok: nivel prerregistrado del contraste
inconclusive. Requiring both tests to pass is an intersection--union test, whose size does not exceed the nominal level.
\end{enumerate}

\subsection{Power, written before estimation}

Table~\ref{tab:power} reports the probabilities of each outcome of the full decision tree, gate included, from 1,000 Monte Carlo
replications with the actual shocks and states, under homoskedastic noise and under noise whose variance rises with $|\varepsilon_t|$.
If the true responses in both states equal the reference effect ($c=1$), the most likely outcome is that the gate fails, and the
probability of confirming H1 is close to the nominal level of the test. The minimum detectable effect (MDE), the $c$ at which the
probability of confirming H1 reaches 0.8, is a $c$ of about five: roughly 10 pp of $\bar y$ per pp of shock. Under the rival hypothesis that the % num-ok: 0.8 es la potencia objetivo que define el MDE
constrained state does not respond ($\beta_C=0$) while the unconstrained state responds with $c\,\beta_{\text{ref}}$, the
probability of confirming H1 stays below the nominal level across the simulated grid. We preregistered H1 as a bounded measurement:
with a reference-sized effect, the design was expected to end at the gate.

\begin{table}[t]
\centering
\caption{Outcome probabilities of the preregistered decision rule, computed before estimation}\label{tab:power}
\begin{tabular}{lcc}
\toprule
 & Homoskedastic & Heteroskedastic \\
\midrule
\multicolumn{3}{l}{\emph{True effect $c=1$ in both states ($\beta_C=\beta_U=\beta_{\text{ref}}$)}} \\
\quad P(gate fails: H0 not confirmed) & \PwHomoHzeroFalla & \PwHetHzeroFalla \\
\quad P(H1 confirmed) & \PwHomoConfirma & \PwHetConfirma \\
\quad P(inconclusive) & \PwHomoNoConc & \PwHetNoConc \\
\quad Power of the H0 test & \PwHomoPowHzero & \PwHetPowHzero \\
\quad Power of the H1 test, full sample & \PwHomoPowHone & \PwHetPowHone \\
\addlinespace
\multicolumn{3}{l}{\emph{Minimum detectable effect (P(H1 confirmed) $=0.8$)}} \\ % num-ok: potencia objetivo del MDE
\quad $c$ & \MdeHomoC & \MdeHetC \\
\quad pp of $\bar y$ per pp of shock & \MdeHomoEps & \MdeHetEps \\
\quad pp of $\bar y$ per pp of federal funds on impact & \MdeHomoFfr & \MdeHetFfr \\
\addlinespace
\multicolumn{3}{l}{\emph{Size and rival}} \\
\quad P(H1 confirmed) with $\beta=0$ & \SizeHomo & \SizeHet \\
\quad Max.\ P(H1 confirmed) under the rival ($\beta_C=0$) & \MdeHomoRivalMax & \MdeHetRivalMax \\
\bottomrule
\end{tabular}

\vspace{0.5em}
\begin{minipage}{\notewidth}\footnotesize
\emph{Notes:} 1,000 replications per cell; actual shocks and states; the estimator, standard errors and decision tree of
Section~\ref{sec:design}. The MDE is interpolated linearly on the simulated grid of $c$; it should be read as roughly 10 pp per
pp of shock and not to more precision. One unit of $c$ is \mdeEpsUnit{} pp of $\bar y$ per pp of shock and \mdeFfrUnit{} pp per pp
of federal funds on impact. The rival maximum is taken over the simulated grid.
\end{minipage}
\end{table}

\subsection{Preregistered robustness}

Four variants of \texttt{s1r\_liq} were preregistered as robustness. They do not enter the decision and are published whatever
their result: \texttt{r1r\_from1971q3} drops 1969Q1--1971Q2, whose shock lags precede 1969 and are set to zero;
\texttt{r1r\_state\_q25} defines the constrained state as the bottom quartile of $s_{t-1}$; \texttt{r1r\_state\_logistic} replaces
$D_{t-1}$ with a logistic function of the standardized $s_{t-1}$, with slope $\gamma=3$ and centered at the median; and
\texttt{r1r\_lags4} uses four lags of the controls.

\section{Results}\label{sec:results}

\subsection{Estimates}

Table~\ref{tab:estimates} reports every preregistered estimate: the linear response, the six estimates of $\beta_C$ (two
confirmatory specifications and four robustness variants) and the six of $\beta_U$. The main specifications are replicated in R
and coincide with the Python estimates to numerical precision. An independent replication by the project's referee, written from
the specification without the authors' code, reproduces all coefficients, the diagnostics and the outcome.

\begin{table}[t]
\centering
\caption{Preregistered estimates: average response of investment, $h=4,\dots,12$}\label{tab:estimates}
\begin{tabular}{llrrrrc}
\toprule
Specification & Term & Estimate & s.e. & d.f. & $p$ (one-sided) & 95\% CI \\ % num-ok: nivel del intervalo
\midrule
\multicolumn{7}{l}{\emph{Gate (decides)}} \\
\texttt{s1r\_lin} & $\beta_{\text{lin}}$ & \bLinEst & \bLinSe & \bLinDof & \bLinP & $[\bLinCiLo, \bLinCiHi]$ \\
\addlinespace
\multicolumn{7}{l}{\emph{Confirmatory specifications for H1 (not evaluated: measurement without a verdict)}} \\
\texttt{s1r\_liq} & $\beta_C$ & \bCLiqEst & \bCLiqSe & \bCLiqDof & \bCLiqP & $[\bCLiqCiLo, \bCLiqCiHi]$ \\
 & $\beta_U$ & \bULiqEst & \bULiqSe & \bULiqDof & \bULiqP & $[\bULiqCiLo, \bULiqCiHi]$ \\
\texttt{s1r\_liq\_nov} & $\beta_C$ & \bCLiqNovEst & \bCLiqNovSe & \bCLiqNovDof & \bCLiqNovP & $[\bCLiqNovCiLo, \bCLiqNovCiHi]$ \\
 & $\beta_U$ & \bULiqNovEst & \bULiqNovSe & \bULiqNovDof & \bULiqNovP & $[\bULiqNovCiLo, \bULiqNovCiHi]$ \\
\addlinespace
\multicolumn{7}{l}{\emph{Preregistered robustness (never decides)}} \\
\texttt{r1r\_from1971q3} & $\beta_C$ & \bCRobFromEst & \bCRobFromSe & \bCRobFromDof & \bCRobFromP & $[\bCRobFromCiLo, \bCRobFromCiHi]$ \\
 & $\beta_U$ & \bURobFromEst & \bURobFromSe & \bURobFromDof & \bURobFromP & $[\bURobFromCiLo, \bURobFromCiHi]$ \\
\texttt{r1r\_state\_q25} & $\beta_C$ & \bCRobQuartEst & \bCRobQuartSe & \bCRobQuartDof & \bCRobQuartP & $[\bCRobQuartCiLo, \bCRobQuartCiHi]$ \\
 & $\beta_U$ & \bURobQuartEst & \bURobQuartSe & \bURobQuartDof & \bURobQuartP & $[\bURobQuartCiLo, \bURobQuartCiHi]$ \\
\texttt{r1r\_state\_logistic} & $\beta_C$ & \bCRobLogitEst & \bCRobLogitSe & \bCRobLogitDof & \bCRobLogitP & $[\bCRobLogitCiLo, \bCRobLogitCiHi]$ \\
 & $\beta_U$ & \bURobLogitEst & \bURobLogitSe & \bURobLogitDof & \bURobLogitP & $[\bURobLogitCiLo, \bURobLogitCiHi]$ \\
\texttt{r1r\_lags4} & $\beta_C$ & \bCRobLagsEst & \bCRobLagsSe & \bCRobLagsDof & \bCRobLagsP & $[\bCRobLagsCiLo, \bCRobLagsCiHi]$ \\
 & $\beta_U$ & \bURobLagsEst & \bURobLagsSe & \bURobLagsDof & \bURobLagsP & $[\bURobLagsCiLo, \bURobLagsCiHi]$ \\
\bottomrule
\end{tabular}

\vspace{0.5em}
\begin{minipage}{\notewidth}\footnotesize
\emph{Notes:} OLS; dependent variable $\bar y_t$, in pp of average investment over $h=4,\dots,12$; coefficients in pp of $\bar y$
per pp of shock (shock summed over the quarter). Shock lags $t-1,\dots,t-10$ and leads $t+1,\dots,t+12$ as controls, not interacted
with the state. Newey--West (Bartlett, five lags) on HC3-rescaled residuals; $t$ on $N-k$ d.f. The $p$-value is one-sided for
$\beta<0$; the interval is two-sided, so the two objects have different levels. Outcome of the preregistered rule:
\texttt{\verdictoVs}. Because the gate fails, no row other than $\beta_{\text{lin}}$ is evaluated: the remaining $p$-values are
descriptive. Source: \texttt{code/analysis/01\_s1r\_estimate.py}; \texttt{output/tables/s1r\_estimates.md}.
\end{minipage}
\end{table}

\subsection{The preregistered outcome, step by step}

\begin{enumerate}
\item \textbf{Gate.} The linear response is $\hat\beta_{\text{lin}}=\bLinEst$ (s.e.\ \bLinSe, \bLinDof{} d.f.), with one-sided
$p=\bLinP$. The gate requires $p<0.05$, so H0 is not confirmed. % num-ok: nivel prerregistrado del contraste
\item \textbf{Step 2 is not evaluated.} Under the rule, $\hat\beta_C$ and $\hat\beta_U$ of \texttt{s1r\_liq} are published as
measurement, without a verdict.
\item \textbf{Step 3 is not evaluated.} The estimates of \texttt{s1r\_liq\_nov} are published, as the preregistration requires,
also when they do not decide.
\end{enumerate}

H0 is not confirmed and H1 is not interpreted: it is neither confirmed, nor refuted, nor inconclusive, because the rule does not
reach it. The six estimates of $\beta_C$ in Table~\ref{tab:estimates}, and their one-sided $p$-values, carry no inferential
weight under the preregistered rule. All six rows are reported, whatever their $p$-values, so that readers see the whole
measurement rather than a selection of it.

\subsection{Shock diagnostics}

Table~\ref{tab:diagnostics} reports the shock diagnostics. The predictability test does not reject at the 5 percent level in
either sample, and the balance test does not reject that the shock is unrelated to the state. The Ljung--Box test does not
reject in the full sample, but it does reject without the Volcker quarters ($p=\DiagNoVLBP$). This was known before estimation,
because the diagnostics do not use investment. As fixed in advance, it changes neither the specification nor the outcome. Its
consequence is that, in \texttt{s1r\_liq\_nov}, the isolated-shock estimate is not comparable one for one with the reference effect,
which is a total response, and a reading in multiples of $\beta_{\text{ref}}$ is only approximate; the relevant scale there is
the MDE in pp of $\bar y$ per pp of shock. Since \texttt{s1r\_liq\_nov} does not decide, the caveat affects only how that
measurement is read.

\begin{table}[t]
\centering
\caption{Shock diagnostics (no investment data used)}\label{tab:diagnostics}
\begin{tabular}{lcc}
\toprule
 & Full sample & Without 1979Q4--1982Q4 \\
\midrule
Quarters ($N$) & \DiagFullN & \DiagNoVN \\
Quarters, constrained / unconstrained state & \DiagFullDoneRestr{} / \DiagFullDoneHolg & \DiagNoVDoneRestr{} / \DiagNoVDoneHolg \\
$\sum|\varepsilon_t|$, constrained / unconstrained & \DiagFullDtwoRestr{} / \DiagFullDtwoHolg & \DiagNoVDtwoRestr{} / \DiagNoVDtwoHolg \\
Balance: $F$ ($p$) & \DiagFullDthreeF{} (\DiagFullDthreeP) & \DiagNoVDthreeF{} (\DiagNoVDthreeP) \\
Ljung--Box(12) ($p$) & \DiagFullLB{} (\DiagFullLBP) & \DiagNoVLB{} (\DiagNoVLBP) \\
Predictability: $F$(\DiagFullPredDfA, \DiagFullPredDfB) ($p$) & \DiagFullPredF{} (\DiagFullPredP) & \\
Predictability: $F$(\DiagNoVPredDfA, \DiagNoVPredDfB) ($p$) & & \DiagNoVPredF{} (\DiagNoVPredP) \\
\bottomrule
\end{tabular}

\vspace{0.5em}
\begin{minipage}{\notewidth}\footnotesize
\emph{Notes:} State $D_{t-1}$ defined with the full-sample median of $s_{t-1}$. Balance: OLS of $\varepsilon_t$ on a constant,
$D_{t-1}$ and $s_{t-1}$, HC1 Wald $F$ on the two slopes. Predictability: OLS of $\varepsilon_t$ on a constant and four lags of
real GDP growth, inflation and the change in the federal funds rate, HC1 Wald $F$ on the twelve slopes. The Ljung--Box statistic
without Volcker concatenates 1979Q3 and 1983Q1. Source: \texttt{output/tables/s1r\_shock\_diagnostics.md}.
\end{minipage}
\end{table}

\section{What the measurement can and cannot rule out}\label{sec:power}

The 95 percent interval for the linear response, $[\bLinCiLo, \bLinCiHi]$, contains zero and the reference effect
$\beta_{\text{ref}}=\betaRef$. The data therefore do not distinguish between no response and a response of the size of the
reference effect. What the interval does exclude, at that level, are average responses more negative than \bLinCiLo{} pp of
$\bar y$ per pp of shock and positive responses above \bLinCiHi.

The outcome is the one the design anticipated. With a reference-sized effect in both states, the probability that the gate
fails was \PwHomoHzeroFalla{} (\PwHetHzeroFalla{} with heteroskedastic noise), and the probability of confirming H1 was
\PwHomoConfirma{} (\PwHetConfirma). The design confirms H1 with probability 0.8 only if the responses in both states are of roughly % num-ok: potencia objetivo del MDE
10 pp of $\bar y$ per pp of shock. Effects of the size of the reference effect are well below what this design can detect with
U.S. aggregate data over this period. The design relies on local projections, which have substantially higher variance than
VAR estimators at medium and long horizons \citep{LiPlagborgMollerWolf2024}; its statistic averages horizons four to twelve, and
its sample has \NShockFull{} quarters split into two states.

The measurement says nothing about state dependence. It neither supports nor contradicts the reserve margin of
Section~\ref{sec:model}.

\section{Limitations}\label{sec:limitations}

\begin{itemize}
\item \textbf{Power.} The design was preregistered as a bounded measurement. Its MDE is about five times the reference effect.
\item \textbf{Estimand.} The coefficients measure the response to an isolated shock. They equal the total response only if the
shocks are serially uncorrelated in the population. Without Volcker the Ljung--Box test rejects, so the comparison with the
reference effect is approximate in that sample.
\item \textbf{Reference effect.} The published figure we translate does not state its horizon, and the ramp that spreads it over
horizons four to twelve is our assumption.
\item \textbf{State.} Aggregate liquidity relative to trend is a proxy for the mass of constrained firms; the model does not
deliver that mapping. The state is predetermined but not exogenous, and differences between states are not effects of liquidity.
\item \textbf{Inference.} The standard errors tend to overcover in the Monte Carlo, which lowers power.
\item \textbf{Data.} Series are in their latest vintage, not real time. The Romer--Romer extension beyond 1996 relies on a
re-estimation with updated Greenbook data.
\item \textbf{Theory.} The propositions are local and partial equilibrium, and they hold the mass of constrained firms fixed. They
assume non-negative net exposure to repricing ($g\ge0$).
\end{itemize}

\section{Conclusion}\label{sec:conclusion}

The constrained firm of rigid models is a corner. With an adjustable precautionary reserve, its investment response is a weighted
average of the unconstrained response and the repricing channel, with a weight bounded away from zero. A common bank dampens the
aggregate response without reversing it. A preregistered measurement on U.S. data for 1969--2005 cannot confirm that the aggregate
response is negative, and therefore, under its own rule, does not interpret the response in the low-liquidity state. With the
power computed before estimation, this was the expected outcome for an effect of the size found in the literature.

Two extensions could add power and were set up with gates fixed before looking at any outcome. A cross-section of industries by
firm size would test the exposure gradient of Corollary~\ref{cor:gradient} with time fixed effects; it did not pass its gate,
because the available survey data do not provide enough industry--size groups with continuous series over the required window. A
panel of euro-area countries would interact the shock with contractual exposure to repricing; it did not pass its gate, because
the euro-area interest rate statistics do not publish, by country, lending volumes by initial rate-fixation period. Both remain
future work, conditional on data that meet those requirements.

\bibliographystyle{p01}
\bibliography{references}

\appendix
\section{Symbolic verification}\label{app:verification}

Every result in Section~\ref{sec:model} is verified with SymPy in two scripts in \texttt{theory/}, run with
\texttt{uv run pytest papers/p01-liquidity/theory}. Each positive test has a negative counterpart that must fail.

\begin{center}\footnotesize\setlength{\tabcolsep}{3pt}
\begin{tabular}{ll}
\toprule
Result & Test \\
\midrule
\multicolumn{2}{l}{\emph{\texttt{theory/check\_nesting.py}}} \\
Prop.~\ref{prop:transmission}, $I_u$ & \texttt{test\_transmission\_formula\_from\_kkt}, \texttt{test\_weighted\_form} \\
Prop.~\ref{prop:transmission}, $m_u$, $\lambda_u$ & \texttt{test\_buffer\_and\_multiplier\_formulas} \\
Prop.~\ref{prop:limits}(i) & \texttt{test\_unconstrained\_limit} \\
Prop.~\ref{prop:limits}(ii) & \texttt{test\_rigid\_limit\_B\_to\_infinity\_is\_g}, \\
 & \texttt{test\_tfg1\_minimum\_branches}, \\
 & \texttt{test\_iopo\_2018\_eq7\_is\_rigid\_limit\_with\_full\_repricing} \\
Prop.~\ref{prop:limits}(iii) & \texttt{test\_Q\_is\_strictly\_decreasing\_onto\_open\_unit\_interval}, \\
 & \texttt{test\_B\_is\_in\_open\_interval\_zero\_kappa}, \texttt{test\_weight\_is\_bounded\_below\_by\_w\_min}, \\
 & \texttt{test\_fixed\_buffer\_corner\_gives\_g\_exactly} \\
Cor.~\ref{cor:gradient} & \texttt{test\_dI\_dg\_is\_one\_minus\_w\_and\_its\_sign} \\
\addlinespace
\multicolumn{2}{l}{\emph{\texttt{theory/check\_aggregation.py}}} \\
Prop.~\ref{prop:aggregation}, \eqref{eq:Lu} & \texttt{test\_general\_formula\_any\_number\_of\_types} \\
Prop.~\ref{prop:aggregation}(i) & \texttt{test\_response\_is\_smaller\_than\_without\_common\_bank} \\
Prop.~\ref{prop:aggregation}(ii) & \texttt{test\_representative\_agent\_matches\_tfg2} \\
Prop.~\ref{prop:aggregation}(iii) & \texttt{test\_L\_u\_positive\_with\_nonnegative\_g\_and\_w\_above\_min} \\
\eqref{eq:Lu} vs.\ nonlinear equilibrium & \texttt{test\_general\_formula\_against\_nonlinear\_equilibrium} \\
\bottomrule
\end{tabular}
\end{center}

\todo{The remark after Corollary~\ref{cor:gradient} combines \eqref{eq:Iu} and \eqref{eq:wmin}; ask Theory for a dedicated test of
$I_u>\omega_{\min}p_f/a$ with $g\ge0$, or drop the remark.}

\end{document}
